\documentclass{article}
\usepackage[utf8]{inputenc}
\usepackage[T1]{fontenc}
\usepackage{geometry}
\geometry{
    a4paper,
    left=2.5cm,
    right=2.5cm,
    top=2.5cm,
    bottom=2.5cm
}
\usepackage{graphicx}
\usepackage{float}
\usepackage{hyperref}
\usepackage[english]{babel}
\usepackage{amsmath}
\usepackage{amssymb}
\usepackage{array}
\usepackage{booktabs}
\usepackage{caption}
\usepackage{subcaption}
\usepackage{enumitem}
\usepackage{fancyhdr}
\usepackage{titlesec}
\usepackage{url}
\usepackage{textcomp}
\usepackage{xcolor}
\usepackage{listings}
\usepackage{multirow}
\usepackage{longtable}
\usepackage{tabularx}
\usepackage{times}

% Set up page style
\pagestyle{fancy}
\fancyhf{}
\rfoot{\thepage}
\renewcommand{\headrulewidth}{0pt}

% Title formatting
\titleformat{\chapter}[display]{\normalfont\bfseries}{\chaptertitle\thechapter}{0pt}{\Huge}
\titlespacing*{\chapter}{0pt}{50pt}{40pt}

% Document info
\title{Template Document}
\author{Academic Support Agent}
\date{\today}

\begin{document}
\maketitle

% Paginas iniciais (sem numera o)
\pagestyle{empty}

% P gina de rosto (capa interna)
% Content from pre-textuais/rosto would be included here


% Ficha catalogr fica: os dados catalogr ficos devem ser fornecidos
% pela BCZM.
% S  s o inclu dos na vers o final da tese ou disserta o. N o s o
% inclu dos nem na proposta de tema de qualifica o nem na vers o
% preliminar da tese ou disserta o: nestes casos, comente a pr xima linha.
% Content from pre-textuais/catalograficos would be included here


% Assinaturas da banca, dedicat ria e agradecimentos
% S  s o inclu dos na vers o final da tese ou disserta o. N o s o
% inclu dos nem na proposta de tema de qualifica o nem na vers o
% preliminar da tese ou disserta o: nestes casos, comente a pr xima linha.
% Content from pre-textuais/preambulo would be included here


%
% O espa amento entre linhas (ATEN O A ESTA PARTE!)
%
%%%%%%%%%%%%%%%%%%%%%%%%%%%%%%%%%%%%%%%%%%%%%%%%%%%%%%%%%%%%%%%%%%%%%%%%%%%%
% PARA A VERS O FINAL:
% Deve ser usado espa amento simples nas p ginas de texto
\singlespacing
% PARA A QUALIFICA O E PARA A VERS O INICIAL:
% Deve ser usado espa amento 1 1/2 nas p ginas de texto
%\doublespacing
%%%%%%%%%%%%%%%%%%%%%%%%%%%%%%%%%%%%%%%%%%%%%%%%%%%%%%%%%%%%%%%%%%%%%%%%%%%%

% Resumo/Abstract
% Content from pre-textuais/resumos would be included here


% Paginas introdut rias (com numera o romana)


% Lista de conte do (sum rio, gerado automaticamente)
% Aqui, e em todos os itens antes da introdu o, o comando \phantomsection   utilizado.
% O seu uso   neecss rio para auxiliar o pacote "hyperref" a fazer a refer ncia correta
% dos links do sum rio, das listas (de tabelas, figuras, algoritmos) com as p ginas 
% respectivas.
% Caso seja tirado, o "hyperref" ir  apontar o link do sum rio para o abstract, o link
% do sum rio para a lista de figuras, o link das listas de figuras para a lista de tabelas,
% e assim sucessivamente.
\phantomsection
\addcontentsline{toc}{chapter}{Sum rio}


% Lista de figuras (gerada automaticamente)

\phantomsection
\addcontentsline{toc}{chapter}{Lista de Figuras}


% Lista de tabelas (gerada automaticamente)

\phantomsection
\addcontentsline{toc}{chapter}{Lista de Tabelas}


% Gloss rio (gerado automaticamente - veja entradas em
% tex/00-introducao/introducao.tex e em tex/02-estilo/estilo.tex)

\phantomsection
\renewcommand{\nomname}{Lista de S mbolos e Abreviaturas}
\markboth{\MakeUppercase{\nomname}}{\MakeUppercase{\nomname}}
\addcontentsline{toc}{chapter}{\nomname}
% O argumento opcional do comando \printnomenclature   a largura deixada
% para os s mbolos no gloss rio. Se seus s mbolos s o "largos", como
% neste exemplo,    melhor por mais espa o do que o 1cm que   reservado
% por default
\printnomenclature[1.5cm]

% P ginas do texto principal (com cabe alho)

\pagestyle{headings}

% Para facilitar a organiza o, foi criado um diret rio para cada
% cap tulo do documento, pois assim os arquivos das figuras ficam
% classificados por cap tulos

% Cap. 1 - Introdu o
% Content from textuais/introducao would be included here


% Cap. 2 - Teoria (Referencial Te rico)
% Content from textuais/teoria would be included here


% Cap. 3 - Trabalhos Relacionados
% Content from textuais/trabalhos_relacionados would be included here


% Cap. 4 - Problema
% Content from textuais/problema would be included here


% Cap. 5 - Implementa o
% Content from textuais/implementacao would be included here


% Cap. 6 - Experimentos e Resultados
% Content from textuais/experimentos_e_resultados would be included here


% Cap. 7 - Conclus o
% Content from textuais/conclusao would be included here


% Refer ncias bibliog ficas (geradas automaticamente)
% Aqui, o comando \phantomsection   utilizado para auxiliar o pacote "hyperref" a fazer a
% refer ncia correta dos links das refer ncias bibliogr ficas com a p gina respectiva.
% Caso seja tirado, o "hyperref" ir  apontar o link das refer ncias bibliogr ficas para a
%  ltima subse o da conclus o.
\phantomsection
\addcontentsline{toc}{chapter}{Refer ncias bibliogr ficas}




%Ap ndice A
% Content from textuais/apendice would be included here
\end{document}